\documentclass[11pt,a4paper]{article}

\usepackage[a4paper,margin=1in]{geometry}
\usepackage{iftex}
\ifPDFTeX
  \usepackage[T2A]{fontenc}
  \usepackage[utf8]{inputenc}
  \usepackage[russian,english]{babel}
\else
  \usepackage{fontspec}
  \usepackage{polyglossia}
  \setdefaultlanguage{russian}
  \setotherlanguage{english}
  \setmainfont{DejaVu Serif}
  \setsansfont{DejaVu Sans}
  \setmonofont{DejaVu Sans Mono}
\fi
\usepackage{amsmath,amssymb,amsthm,mathtools}
\usepackage{hyperref}
\usepackage{enumitem}
\usepackage{xcolor}
\usepackage{booktabs}
\usepackage{longtable}
\usepackage{array}

\hypersetup{
  colorlinks=true,
  linkcolor=blue!50!black,
  urlcolor=blue!60!black,
  citecolor=blue!50!black
}

\title{Промежуточный прогресс по Multi-Agent Nash Learning from HF:\\
теория, реализация в VERL, эксперименты и дальнейший план}
\author{}
\date{}

\begin{document}
\maketitle

\section{Постановка задачи}

Проект находится на пересечении \textit{preference learning}, \textit{game dynamics} и \textit{multi-agent LLM training}. Базовая мотивация идёт от статьи \textit{Nash Learning from Human Feedback}~\cite{nlhf}, где предлагается вместо явной reward-модели работать с более общей моделью предпочтений. Это важно, потому что человеческие предпочтения в общем случае могут быть \textbf{нетранзитивны}: может выполняться
\[
A \succ B, \qquad B \succ C, \qquad C \succ A.
\]
В таком случае сведение качества ответа к одной скалярной награде может быть слишком грубым или даже концептуально неверным.

Основная идея проекта состоит в том, чтобы перейти от постановки ``модель играет против прошлых версий себя'' к постановке, где:
\begin{enumerate}[label=\arabic*.]
  \item либо играют \textbf{две разные модели} с разными pretrain / tokenizer / inductive bias;
  \item либо играют \textbf{две версии одной и той же базы}, но обученные на разных задачах и потому обладающие разными навыками.
\end{enumerate}

Цель --- понять, можно ли получить более интересную и потенциально более устойчивую динамику обучения, чем в обычном RLHF, если обучать модели в игровом фреймворке и обновлять их на основе попарных сравнений ответов.

\section{Почему обычного RLHF недостаточно}

Классический RLHF обычно записывается как задача
\[
\max_{\pi_\theta}
\;
\mathbb{E}_{x \sim \mathcal D,\; y \sim \pi_\theta(\cdot \mid x)}
\bigl[r_\phi(x,y)\bigr]
- \beta \, \mathrm{KL}(\pi_\theta \,\|\, \pi_{\mathrm{ref}}),
\]
где:
\begin{itemize}
  \item $x$ --- prompt;
  \item $y$ --- ответ модели;
  \item $r_\phi(x,y)$ --- reward model;
  \item $\pi_{\mathrm{ref}}$ --- reference policy;
  \item $\beta$ --- коэффициент регуляризации.
\end{itemize}

Такая формулировка предполагает, что качество ответа можно описать единственным числом $r_\phi(x,y)$. Однако в задачах с человеческими предпочтениями естественнее иметь объект вида
\[
P_\phi(y_a \succ y_b \mid x),
\]
то есть вероятность того, что при prompt $x$ ответ $y_a$ предпочтительнее ответа $y_b$.

Если бы существовала ``истинная'' reward-функция, то часто можно было бы писать
\[
P(y_a \succ y_b \mid x) \approx \sigma\bigl(r(x,y_a)-r(x,y_b)\bigr),
\]
но в общем случае pairwise preference может быть богаче, чем любая одна скалярная reward-функция. Именно здесь появляется мотивация для \textit{Nash-style} постановки: искать не модель с максимальным reward в вакууме, а стратегию, устойчивую к другим стратегиям.

\section{Интуиция Nash-постановки}

Если рассматривать две политики $\pi_A$ и $\pi_B$ как игроков, то естественно ввести utility-функцию $u(\pi_A,\pi_B)$, описывающую, насколько хорошо играет $A$ против $B$. Тогда в двухигроковой постановке равновесие Нэша $(\pi_A^\star,\pi_B^\star)$ удовлетворяет
\[
u(\pi_A^\star,\pi_B^\star) \ge u(\pi_A,\pi_B^\star)
\quad \forall \pi_A,
\]
\[
u(\pi_A^\star,\pi_B^\star) \le u(\pi_A^\star,\pi_B)
\quad \forall \pi_B.
\]

Даже если текущая реализация не является строгой нулевой суммой и не даёт формального доказательства достижения равновесия Нэша, сама идея проекта уже строится в правильном направлении: \textbf{оценивать качество модели через игру с альтернативными стратегиями, а не только через скалярный reward}.

\section{Два предполагаемых исследовательских сетапа}

\subsection{Сетап 1: разные производители, разные модели}

В первом сетапе играют две модели разных семейств, например:
\begin{itemize}
  \item \texttt{Llama} против \texttt{Qwen};
  \item или, как в текущем эксперименте, \texttt{Qwen} против \texttt{SmolLM}.
\end{itemize}

Гипотеза здесь такая: разные модели обладают разными знаниями, разными токенизаторами, разными inductive bias и потому будут систематически выигрывать на разных подмножествах задач. Если затем обновлять каждую модель на тех задачах, где она проиграла сопернику, можно получить targeted improvement на собственных слабых местах.

\subsection{Сетап 2: одна база, разные специализации}

Во втором сетапе рассматриваются две версии одной и той же базовой модели, но затюненные на разные типы задач. Например:
\begin{itemize}
  \item одна лучше в math;
  \item другая лучше в code;
  \item или одна сильнее в reasoning, а другая --- в instruction-following.
\end{itemize}

Гипотеза: через игровую динамику такие модели могут частично перенять навыки друг друга. Это потенциально интереснее с научной точки зрения, чем первый сетап, потому что здесь меньше архитектурного шума: модели одной семьи проще сравнивать и проще интерпретировать возникающую динамику как эффект preference/game learning, а не как побочный эффект разных backbone.

\section{Роль VERL в этом проекте}

В данном проекте \texttt{VERL} играет роль не отдельного алгоритма, а \textbf{исследовательской инфраструктуры}. Он даёт:
\begin{itemize}
  \item распределённую архитектуру через \texttt{Ray};
  \item абстракцию \texttt{DataProto} для данных и meta-information;
  \item роли \texttt{Actor}, \texttt{Rollout}, \texttt{RefPolicy};
  \item FSDP и related distributed training mechanics;
  \item recipe-уровень для быстрого прототипирования новых training loops.
\end{itemize}

С инженерной точки зрения задача проекта не сводится к написанию одного loss. Нужно уметь:
\begin{itemize}
  \item генерировать ответы несколькими политиками;
  \item сравнивать их через внешний judge;
  \item собирать pairwise update data;
  \item считать reference log-probabilities;
  \item сохранять и восстанавливать checkpoint'ы;
  \item валидировать начальные и обученные веса в нескольких матчапах.
\end{itemize}

Именно поэтому построение working recipe в \texttt{VERL} само по себе уже является существенным исследовательским прогрессом.

\section{Архитектура репозитория по смыслу}

На содержательном уровне текущий код разбивается на следующие слои:
\begin{itemize}
  \item \texttt{recipe/spin} --- базовый online-DPO/self-play стек;
  \item \texttt{recipe/refplay} --- более прикладной стек обучения от сравнений ответов;
  \item \texttt{recipe/nlhf} --- baseline, близкий к исходной статье \textit{Nash Learning from Human Feedback};
  \item \texttt{recipe/crossplay} --- новый рецепт для конкурентного обучения двух разных живых моделей.
\end{itemize}

Это удобно понимать как лестницу:
\begin{enumerate}[label=\arabic*.]
  \item \texttt{SPIN} даёт общую механику online preference/DPO.
  \item \texttt{RefPlay} даёт обвязку для training from comparisons.
  \item \texttt{NLHF} реализует baseline в духе paper setup.
  \item \texttt{CrossPlay} --- это собственное расширение в сторону multi-model game.
\end{enumerate}

\section{Что реализовано в paper-faithful NLHF baseline}

\texttt{NLHF} в текущем коде реализован как dual-policy recipe поверх \texttt{refplay}-стека. Схема следующая:
\begin{enumerate}[label=\arabic*.]
  \item \texttt{policy\_a} генерирует ответ на prompt;
  \item \texttt{policy\_b} генерирует ответ на тот же prompt;
  \item pairwise judge вычисляет preference probability;
  \item frozen reference policy даёт reference log-probabilities;
  \item trainable policy обновляется с учётом предпочтения и KL-штрафа.
\end{enumerate}

Пусть $p \in [0,1]$ --- вероятность того, что judge предпочитает ответ текущей policy. Тогда в коде используется величина
\[
A_{\mathrm{pref}} = p - 0.5,
\]
которая играет роль ``preference advantage''. Далее оценивается KL-like отклонение от reference:
\[
\widehat{\mathrm{KL}} = \log \pi_\theta(y \mid x) - \log \pi_{\mathrm{ref}}(y \mid x),
\]
после чего строится вес:
\[
w = A_{\mathrm{pref}} - \tau \widehat{\mathrm{KL}},
\]
и оптимизируется loss
\[
\mathcal L_{\mathrm{NLHF}}
=
-
\mathbb E\bigl[\mathrm{stopgrad}(w)\cdot \log \pi_\theta(y \mid x)\bigr].
\]

Интуитивно:
\begin{itemize}
  \item если judge предпочитает ответ модели, то $p > 0.5$, и policy получает положительный learning signal;
  \item если ответ хуже альтернативного, то $p < 0.5$, и learning signal становится отрицательным;
  \item коэффициент $\tau$ регулирует силу удержания относительно reference.
\end{itemize}

Такой baseline уже ценен как paper-faithful точка отсчёта. Однако он всё ещё ближе к схеме ``одна главная policy против alternative/opponent, синхронизируемого через EMA'', чем к полностью симметричной игре двух независимых агентов.

\section{Что реализовано в CrossPlay}

\texttt{CrossPlay} --- это основной практический вклад текущего этапа работы. Его ключевая идея: \textbf{обе модели живые, обе генерируют ответы на одинаковые prompt'ы, и каждая обучается только на собственных поражениях}.

\subsection{Высокоуровневый пайплайн}

Для каждого prompt $x$:
\[
y_A \sim \pi_A(\cdot \mid x),
\qquad
y_B \sim \pi_B(\cdot \mid x).
\]

Judge независимо оценивает оба ответа:
\[
s_A = J(x, y_A),
\qquad
s_B = J(x, y_B).
\]

Дальше:
\begin{itemize}
  \item если $s_B > s_A$, то \texttt{policy\_a} проиграла;
  \item для \texttt{policy\_a} создаётся пара
  \[
  y^+ = y_B,\qquad y^- = y_A;
  \]
  \item если $s_A > s_B$, то симметрично проиграла \texttt{policy\_b}.
\end{itemize}

То есть каждая policy получает update не на каком-то абстрактном ``хорошем'' ответе, а на том конкретном ответе соперника, которым она была побита на данной задаче.

\subsection{Почему это не просто ещё один DPO}

Здесь есть три принципиально важных расширения.

\paragraph{1. Гетерогенность моделей.}
\texttt{policy\_a} и \texttt{policy\_b} могут быть из разных model families и иметь разные токенизаторы. Это значит, что обучение строится не на shared token space, а на текстовом сравнении:
\begin{enumerate}[label=\arabic*.]
  \item обе модели генерируют текст;
  \item judge судит текст;
  \item chosen/rejected пары потом отдельно ретокенизируются токенизатором целевой policy.
\end{enumerate}

\paragraph{2. Раздельные anchors.}
У каждой модели свой frozen anchor:
\begin{itemize}
  \item \texttt{anchor\_a} для \texttt{policy\_a};
  \item \texttt{anchor\_b} для \texttt{policy\_b}.
\end{itemize}
Это особенно важно для cross-vendor матчапов.

\paragraph{3. Benchmark-aware replay.}
Ошибки не просто складываются в один общий буфер. Они группируются по benchmark buckets, и sampling при update усиливает именно те подмножества задач, где политика систематически проигрывает.

\section{Математика CrossPlay-обновления}

Пусть \texttt{policy\_a} проиграла на prompt $x$. Тогда для неё строится DPO-пара:
\[
y^+ = y_B,\qquad y^- = y_A.
\]

Соответствующий reference-DPO loss можно записать как
\[
\mathcal L_{\mathrm{DPO}}^A
=
-
\log \sigma \Bigl(
\beta
\bigl[
(\log \pi_A(y^+ \mid x)-\log \pi_A(y^- \mid x))
-
(\log \rho_A(y^+ \mid x)-\log \rho_A(y^- \mid x))
\bigr]
\Bigr),
\]
где:
\begin{itemize}
  \item $\pi_A$ --- текущая \texttt{policy\_a};
  \item $\rho_A$ --- её frozen anchor;
  \item $\beta$ --- DPO temperature.
\end{itemize}

Аналогично строится loss для \texttt{policy\_b}, если она проиграла \texttt{policy\_a}.

Таким образом, \texttt{CrossPlay} по духу сочетает:
\begin{itemize}
  \item game-based data generation;
  \item pairwise preference construction;
  \item reference-regularized policy updates.
\end{itemize}

\section{Benchmark-aware replay более формально}

Для каждой policy хранится replay-buffer её собственных неудач. Для benchmark bucket $b$ вычисляется priority, которая интуитивно растёт, если:
\begin{itemize}
  \item политика чаще проигрывает на этом benchmark'е;
  \item и если отставание по среднему reward там больше.
\end{itemize}

Схематично это выглядит как
\[
\mathrm{priority}(b)
=
1
+
\max(\mathrm{opp\_wins}_b-\mathrm{own\_wins}_b,0)
+
\max(\mathrm{reward\_gap}_b,0).
\]

Внутри benchmark bucket отдельные примеры тоже могут перевзвешиваться по margin:
\[
w_{\mathrm{example}} \propto (1+\mathrm{margin})^\alpha.
\]

Это очень важный момент. `CrossPlay` не просто обучает каждую модель на ``любых поражениях'', а создаёт механизм \textbf{фокусированного подтягивания слабых областей}.

\section{Какой сетап уже реально реализован и проверен}

На текущий момент реально реализован, запущен и доведён до конца heterogeneous \texttt{CrossPlay} setup:
\begin{itemize}
  \item \texttt{policy\_a}: \texttt{Qwen/Qwen2.5-0.5B-Instruct};
  \item \texttt{policy\_b}: \texttt{HuggingFaceTB/SmolLM2-1.7B-Instruct}.
\end{itemize}

Датасет:
\begin{itemize}
  \item train: \texttt{gsm8k/train\_plain.parquet};
  \item val: \texttt{gsm8k/test\_plain.parquet}.
\end{itemize}

Размеры:
\begin{itemize}
  \item train: $7473$ примера;
  \item val: $1319$ примеров.
\end{itemize}

Основные training hyperparameters:
\begin{itemize}
  \item \texttt{train\_batch\_size} $= 1$;
  \item \texttt{ppo\_mini\_batch\_size} $= 1$;
  \item \texttt{ppo\_micro\_batch\_size\_per\_gpu} $= 1$;
  \item \texttt{val\_batch\_size} $= 32$;
  \item \texttt{total\_training\_steps} $= 3737$.
\end{itemize}

С инженерной точки зрения важно, что полный сценарий был:
\begin{enumerate}[label=\arabic*.]
  \item запущен;
  \item упал по памяти на промежуточном шаге;
  \item безопасно resumed;
  \item доведён до финального checkpoint'а.
\end{enumerate}

Это существенный результат: в distributed RLHF-пайплайнах нередко именно reproducible end-to-end execution является главным риском.

\section{Текущий judge и его ограничения}

Текущий judge для \texttt{GSM8K} очень грубый. Он выдаёт только три значения:
\[
1.0 \quad \text{за exact match},
\]
\[
0.1 \quad \text{за ответ, из которого извлеклось число, но оно неверно},
\]
\[
0.0 \quad \text{если final answer не извлёкся}.
\]

Причём извлечение завязано на формат
\[
\texttt{\#\#\#\# ...}
\]
и последующее вытаскивание числа из области final answer.

Это означает, что множество качественно разных ошибок схлопывается в одну и ту же reward-категорию. С точки зрения signal quality это главный текущий bottleneck.

\section{Экспериментальные результаты}

На eval по $500$ validation examples были получены следующие outcome rates:

\begin{center}
\begin{tabular}{lccccc}
\toprule
Run & $A$ win & $B$ win & Tie & $\mathbb E[r_A]$ & $\mathbb E[r_B]$ \\
\midrule
\texttt{init\_vs\_init} & 0.008 & 0.084 & 0.908 & 0.0026 & 0.0084 \\
\texttt{initA\_vs\_ckptB} & 0.008 & 0.492 & 0.500 & 0.0026 & 0.0548 \\
\texttt{ckpt\_vs\_ckpt} & 0.506 & 0.006 & 0.488 & 0.1000 & 0.0548 \\
\texttt{ckptA\_vs\_initB} & 0.916 & 0.000 & 0.084 & 0.1000 & 0.0084 \\
\bottomrule
\end{tabular}
\end{center}

\subsection{Как это интерпретировать}

Из этих чисел следует:
\begin{enumerate}[label=\arabic*.]
  \item На старте \texttt{SmolLM 1.7B} сильнее \texttt{Qwen 0.5B}.
  \item После обучения выросли обе стороны.
  \item Наиболее сильный рост по текущему judge произошёл у \texttt{policy\_a}.
\end{enumerate}

Однако крайне важно не переинтерпретировать эти результаты. Если посмотреть на reward means, то значение порядка $0.1$ соответствует не ``модель хорошо решает задачи'', а скорее режиму ``модель стабильно выдаёт parseable final answer, но далеко не всегда правильный''. Поэтому текущий эксперимент показывает:
\begin{itemize}
  \item \textbf{не} окончательную победу в math reasoning;
  \item а то, что competitive training loop уже даёт measurable effect under current judge.
\end{itemize}

\section{Главный научный вывод промежуточного этапа}

Главный вывод текущего этапа можно сформулировать так:

\begin{quote}
\textbf{Основная инженерная гипотеза подтверждена: heterogeneous competitive preference training в VERL можно реализовать, стабилизировать и довести до end-to-end эксперимента. Главный bottleneck следующего этапа находится уже не в training loop, а в качестве reward/preference signal.}
\end{quote}

Это очень хороший промежуточный результат. Он означает, что проект уже перешёл из фазы ``сможем ли мы вообще запустить такую постановку'' в фазу ``какой judge / reward / preference design действительно ведёт к содержательному улучшению''.

\section{Что можно и чего нельзя пока утверждать}

\subsection{Что уже можно утверждать честно}

\begin{itemize}
  \item Построена рабочая исследовательская платформа для multi-agent preference learning в \texttt{VERL}.
  \item Реализован paper-like \texttt{NLHF} baseline.
  \item Реализован новый \texttt{CrossPlay} recipe для двух разных живых моделей.
  \item Проведён полный end-to-end experiment с checkpoint recovery и offline eval.
  \item Получен первый measurable signal в пользу competitive training under current judge.
  \item Локализован следующий научный bottleneck: coarse reward / judge.
\end{itemize}

\subsection{Чего пока утверждать нельзя}

\begin{itemize}
  \item Что уже получено строгое Nash equilibrium learning.
  \item Что модель уже доказуемо стала лучше математически в широком смысле.
  \item Что \texttt{CrossPlay} уже превосходит RLHF или исходный \texttt{NLHF} на честных внешних метриках.
\end{itemize}

\section{Что дальше делать}

Наиболее логичный roadmap выглядит так.

\subsection{Шаг 1. Улучшить сигнал}

Сейчас самая важная задача --- заменить judge $1/0.1/0$ на более выразительный signal. Возможные направления:
\begin{itemize}
  \item richer outcome reward;
  \item pairwise preference model;
  \item verifier-style reward;
  \item process-aware scoring.
\end{itemize}

\subsection{Шаг 2. Сравнить mathematical objectives}

Сейчас в репозитории уже есть как минимум два содержательно разных подхода:
\begin{itemize}
  \item \texttt{NLHF}: preference-probability update;
  \item \texttt{CrossPlay}: competitive data generation + reference DPO.
\end{itemize}

Следующий шаг --- поставить их рядом как baseline'ы и сравнить:
\begin{itemize}
  \item устойчивость обучения;
  \item скорость адаптации;
  \item качество итоговых матчапов;
  \item чувствительность к типу judge.
\end{itemize}

\subsection{Шаг 3. Расширить данные}

\texttt{GSM8K} удобен как debugging dataset, но дальше для содержательного reasoning желательно переходить к более сложным задачам. Важно, однако, делать это \textbf{после} улучшения judge, иначе будет трудно отделить проблемы датасета от проблем reward signal.

\subsection{Шаг 4. Реализовать и запустить второй сетап}

Сетап ``одна база, разные специализации'' очень важен, потому что:
\begin{itemize}
  \item он даёт более чистую интерпретацию результатов;
  \item он ближе к задаче skill transfer;
  \item он может показать, возможно ли через game dynamics получить multi-skill policy лучше, чем через обычный RLHF на смеси задач.
\end{itemize}

\subsection{Шаг 5. Вернуться к более строгой Nash-интерпретации}

После стабилизации signal design имеет смысл переходить к более строгим вопросам:
\begin{itemize}
  \item exploitability;
  \item циклы предпочтений;
  \item population-based opponents;
  \item устойчивость к смене judge.
\end{itemize}

\section{Итог}

Если формулировать общий результат максимально точно, то он звучит так:

\begin{quote}
За текущий этап была не просто написана новая training script-логика, а построена исследовательская платформа в \texttt{VERL} для multi-agent preference learning. В ней уже поддерживаются paper-faithful \texttt{NLHF} baseline и новый heterogeneous \texttt{CrossPlay} setting, который доведён до полноценного end-to-end эксперимента, resume после OOM и offline evaluation. Первый эксперимент показал, что competitive preference-based обучение действительно даёт measurable signal, а главный limiting factor следующего этапа --- это уже не distributed pipeline, а качество reward/preference signal.
\end{quote}

Самая полезная короткая внутренняя формулировка для понимания собственного прогресса такая:

\begin{quote}
Я уже реализовала и проверила платформу для multi-agent preference learning в VERL; сейчас работа перешла из инженерной фазы ``запустить игру между моделями'' в исследовательскую фазу ``какой preference/reward signal действительно ведёт к содержательному улучшению''.
\end{quote}

\begin{thebibliography}{9}

\bibitem{nlhf}
Liu et al.
\newblock \textit{Nash Learning from Human Feedback}.
\newblock \url{https://arxiv.org/abs/2312.00886}

\end{thebibliography}

\end{document}
